A bank may know that a policy works on average and still need to choose
which mechanism to change, how large an intervention to offer and how much
confidence to place in the causal assumptions. These are related decisions,
but they require different targets. We begin with channel attribution and
feasible mediator interventions, then compare natural-effect estimation and
sensitivity methods across implementations. The final dose analysis asks how
to choose an amount rather than merely fit a response curve. Throughout,
changes in assumptions and changes in sampling uncertainty remain distinct.

\section{Commercial mediation software and the meaning of a channel}
\label{exec20:med-scope}
A credit team may know that a servicing offer improves repayment while still not knowing whether the improvement travels through customer understanding, available liquidity or a change in staff behavior. That distinction matters when the proposed investment changes only one channel. An effect attributed to a mediator is not automatically the effect of a feasible intervention on that mediator. Commercial mediation software supplies developed estimation and reporting workflows, but the comparison must begin with the intervention, counterfactual target and observation clock.

The selected Stata manuals and Mplus chapter document concrete capabilities. The StataNow two-mediator manual is a separate product feature from the one-mediator command. The selected Mplus chapter has March 2017 PDF metadata and supports a historical technical profile; it does not establish current feature absence or exact licensed entitlement. Sources acquired after the literature cutoff cannot, through an internal date alone, authenticate their exact bytes on that cutoff date. The available comparison therefore supports scoped design choices while retaining unresolved versions and commercial alternatives.~\citep{exec20_med_stata_one,exec20_med_stata_multiple,exec20_med_mplus_ch17}

\subsection{Four nested means and two valid decompositions}
\label{exec20:med-targets}
Let $A\in\{0,1\}$ indicate an offer, $M(a)$ the mediator that would follow offer $a$, and $Y(a,m)$ the outcome under offer $a$ and mediator setting $m$. For example, $M$ might describe a measured comprehension score and $Y$ a net recovery amount over a fixed subsequent horizon. The intervention on a score needs its own operational interpretation: administering a quiz, improving understanding and changing a recorded answer are distinct actions. Define
\begin{equation}
 \mu_{ab}=E\{Y(a,M(b))\},\qquad a,b\in\{0,1\}.
 \label{exec20:med-means}
\end{equation}
The means with $a\ne b$ combine an outcome under one offer with a mediator under another. Observing the offer, mediator and outcome on many customers does not directly reveal these nested outcomes.

Using labels that state which component is held fixed, the additive total effect has two decompositions:
\begin{align}
 \mu_{11}-\mu_{00}
  &=(\mu_{10}-\mu_{00})+(\mu_{11}-\mu_{10}),\nonumber\\
  &=(\mu_{11}-\mu_{01})+(\mu_{01}-\mu_{00}).
 \label{exec20:med-decompositions}
\end{align}
The first pairs a pure natural direct effect with a total natural indirect effect. The second pairs a total natural direct effect with a pure natural indirect effect. Each identity follows by adding and subtracting its intermediate mean. Mixing the direct term from one line with the indirect term from the other generally fails to recover the total. Stata's one-mediator manual documents these alternatives and includes treatment--mediator interaction by default. Its supported model combinations also matter: a logit outcome cannot be combined with the listed linear or exponential-mean mediator model in this interface.~\citep{exec20_med_stata_one}

Under the manual's stated identification assumptions, a nested mean is represented by averaging an outcome regression under $A=a$ over the mediator distribution under $A=b$ and the baseline population:
\begin{equation}
 \mu_{ab}=\int\!\!\int
 E(Y\mid A=a,M=m,X=x)\,dF(m\mid A=b,X=x)\,dF(x).
 \label{exec20:med-integral}
\end{equation}
Here $X$ contains permitted pretreatment covariates. The assumptions include treatment and mediator ignorability, appropriate support and the absence of a treatment-induced mediator--outcome confounder for this standard formulation. Regression software can calculate the integral under its models; it cannot infer those scientific conditions from a converged fit. Stata's introductory workflow explicitly places question, decomposition and assumptions before estimation and interpretation.~\citep{exec20_med_stata_one,exec20_med_stata_intro}

\subsection{Interaction changes the channel attribution}
\label{exec20:med-interaction}
An authored continuous example makes the labels consequential. Suppose
\begin{align}
 M(a)&=2+3a+U,\nonumber\\
 Y(a,m)&=5+4a+2m+am+V,
 \label{exec20:med-law}
\end{align}
where $U$ and $V$ have mean zero and are independent under the stipulated causal model. The four means are $\mu_{00}=9$, $\mu_{10}=15$, $\mu_{01}=15$ and $\mu_{11}=24$. The total effect is 15. Its pure-direct/total-indirect decomposition is $6+9$, while its total-direct/pure-indirect decomposition is $9+6$. Multiplying the mediator coefficient 3 by the untreated outcome slope 2 gives 6; that is the pure indirect component, not the total indirect component 9. Dropping the interaction changes the scientific model instead of merely simplifying a report.~\citep{exec20_med_reference}

The controlled direct effect at a fixed mediator value $m$ is
\begin{equation}
 E\{Y(1,m)-Y(0,m)\}=4+m.
 \label{exec20:med-controlled}
\end{equation}
It is 4, 6 or 9 at $m=0,2,5$, respectively. A bank specifying an operational mediator setting should request that setting explicitly. If $Y$ is measured in dollars and the full offer costs an incremental USD 10 per customer, the stipulated total-effect margin is USD 5. Charging that cost against one natural component would answer a different question unless a feasible component intervention and its cost were separately justified. This example supplies a finite analytical reference, not a commercial engine benchmark or a measured banking return.

The inspected Mplus chapter distinguishes conventional coefficient-product effects from counterfactually defined effects. Its conventional branch allows multiple mediators and outcomes; categorical mediator or outcome paths can use underlying latent response variables. A product on a latent-response scale is not thereby a difference in repayment probabilities. The chapter separately documents total natural indirect, pure natural direct and total effects for a single mediator, with \texttt{IND} and \texttt{MOD} syntax and explicit controlled-effect mediator values. The appropriate branch and scale should be part of the requested target, not inferred from the word indirect in an output table.~\citep{exec20_med_mplus_ch17}

The selected \texttt{MODEL INDIRECT} command uses delta-method standard errors by default, with documented bootstrap and interval alternatives. The chapter excludes this command for \texttt{TYPE=RANDOM}, the \texttt{CONSTRAINT} option of the \texttt{VARIABLE} command and \texttt{TYPE=EFA}. Those are statements about this inspected command and version profile. They do not prove that every other Mplus modeling route lacks a desired effect. Likewise, a documented interval option does not establish its calibration under shared bank or calendar shocks.

\subsection{Randomizing an offer does not identify every channel}
\label{exec20:med-nonidentification}
Consider two authored worlds with independent fair binary variables $A$ and $U$. The offer $A$ is randomized, and the mediator is
\begin{equation}
 M(a)=a\mathbin{\oplus}U,
 \label{exec20:med-xor}
\end{equation}
where $\oplus$ denotes exclusive OR. In the direct world, $Y(a,m)=a$. In the mediated world, $Y(a,m)=m\mathbin{\oplus}U$. In both worlds the observed mediator is equally likely to be zero or one within each offer arm, and the observed outcome is exactly $Y=A$. The complete observed joint distribution of $(A,M,Y)$ is identical. There is observed mediator overlap, and the offer's total effect is 1 in both worlds.~\citep{exec20_med_reference}

The nested means differ. In the direct world, $Y(1,M(0))=1$, so the pure natural direct effect is 1 and the total natural indirect effect is 0. In the mediated world,
$Y(1,M(0))=U\oplus U=0$, so those components are 0 and 1. The latter world contains unobserved mediator--outcome confounding through $U$, violating the assumptions needed for the ordinary mediation integral. The construction therefore does not contradict software operating under qualified assumptions. It shows why randomizing $A$ and checking the observed support of $M$ are insufficient to justify the natural-effect decomposition.

Even the mediated world's natural indirect effect of 1 is not the effect of setting everyone's mediator to one rather than zero. Under either constant mediator setting, $m\oplus U$ has mean $1/2$, so this controlled mediator-only contrast is zero. At an assumed USD 100 per beneficial outcome and USD 20 intervention cost, offering $A$ has a stipulated incremental margin of USD 80, whereas the mediator-only contrast gives USD $-20$. These refer to different proposed actions. A channel attribution cannot be sold as the return from changing a dashboard score.

An interventional mediator target would instead specify a feasible distributional intervention, for example a draw from a declared reference mediator distribution conditional on eligible baseline information. That differs from retaining each person's own unobserved $M(b)$ inside a nested counterfactual. Its identification conditions and practical implementation need their own derivation, especially when treatment changes mediator--outcome confounders. Section~\ref{sec:exec25-rmed} develops one such interventional mean and the associated restrictions using the retained primary text. The selected commercial documents do not establish that every such target is available. The proposal should preserve unsupported targets rather than relabel a natural effect as interventional.

\subsection{Population averaging, signs and multiple mediators}
\label{exec20:med-population}
Stata's controlled-effect postestimation averages counterfactual predictions over covariates by default; its \texttt{atmeans} option evaluates them at covariate means. In the inspected Mplus moderation syntax, continuous control variables evaluated at zero can instead be evaluated at their means by centering them. A comparison must therefore specify the target population and the reference point before interpreting numerical disagreement.~\citep{exec20_med_stata_post,exec20_med_mplus_ch17}

For an exact nonlinear illustration, let $X$ be equally likely to equal $-1$ or 1 and stipulate
\begin{equation}
 P\{Y=1\mid do(A=a),do(M=m),X=x\}
   =\tfrac15+a\left(\tfrac1{10}+\tfrac{x^2}{10}\right).
 \label{exec20:med-population-law}
\end{equation}
The mediator has no role in this law. The population-average controlled effect is $1/5$, while evaluation at the mean covariate value $E(X)=0$ gives $1/10$. At an assumed value of USD 100 and cost of USD 15, those two targets imply margins of USD 5 and USD $-5$. The example isolates the averaging operation and does not claim that a selected vendor's available parametric family reproduces this law automatically.~\citep{exec20_med_reference}

The proportion mediated is also a derived quantity whose denominator and signs must remain visible. Under the simple stipulated law $M(a)=a$, $Y(a,m)=3a-2m$, the total effect is 1, the pure direct component is 3 and the total indirect component is $-2$. The ratio of indirect to total is $-2$, not a share between zero and one. With total effect approaching zero, such a ratio can become arbitrarily unstable even when the components are meaningful. Stata's postestimation manual normally refuses conflicting effect signs and permits an explicit \texttt{force} override. That override authorizes a computation, not an interpretation as a percentage of benefit.~\citep{exec20_med_stata_post,exec20_med_reference}

The separate StataNow two-mediator manual documents parallel and causally ordered sequential mediators. Its selected syntax supports continuous outcome and mediator variables, with binary or continuous treatment. Interaction switches distinguish treatment--mediator, mediator--mediator and first-mediator/treatment interactions in the second mediator equation. The stated model requires no unmeasured confounding across the treatment, mediator and outcome system, with independent zero-mean errors in the displayed parametric model. More path components create more scientific assumptions and more possible decompositions; they do not reveal a causal ordering from the data alone.~\citep{exec20_med_stata_multiple}

For a bank, understanding and liquidity may be parallel channels, understanding may change liquidity management, or an unobserved hardship shock may affect both. The business description does not select among these graphs. An assistant can draft alternatives, retrieve source-supported mechanisms and request the relevant measurement clocks. It should retain unresolved arrows and refuse to convert an expert's preferred story into verified mediator independence. The resulting graph and estimand should be reviewed before a vendor command is constructed.

\subsection{Sensitivity parameters do not define probability intervals}
\label{exec20:med-sensitivity}
The absence of a measured confounder can motivate a sensitivity analysis, but the parameter and its interpretation must be explicit. Consider the authored structural model
\begin{equation}
 M=aA+u,\qquad Y=cA+bM+e,
 \label{exec20:med-residual-law}
\end{equation}
where $A$ is independent of the errors, $\operatorname{Var}(u)>0$ and
$\rho=\operatorname{Corr}(u,e)$. After removing the linear contribution of $A$, the coefficient of $M$ in an ordinary outcome regression satisfies
\begin{equation}
 b_{\rm OLS}
   =b+\frac{\operatorname{Cov}(u,e)}{\operatorname{Var}(u)}
   =b+\rho\frac{\sigma_e}{\sigma_u}.
 \label{exec20:med-bias}
\end{equation}
This follows directly by projecting $bu+e$ onto $u$. Under the stated additive structural intervention, the indirect effect is $ab$. If $a=2$, $b_{\rm OLS}=3$ and the structural scale ratio $\sigma_e/\sigma_u=2$ is assumed, varying $\rho$ over $[-1/2,1/2]$ produces effects from 4 to 8, with value 6 at $\rho=0$. At an assumed USD 100 per effect unit and USD 500 cost for a separately justified intervention, the corresponding arithmetic margins range from USD $-100$ to USD 300.~\citep{exec20_med_reference}

This is a set of declared sensitivity scenarios, not a 95\% confidence interval or a posterior credible interval. The structural scale ratio is an assumption; it must not be substituted silently by the residual standard deviation from a regression that has already projected part of $e$ onto $u$. Sampling uncertainty within each scenario is a separate calculation. The scenarios have not been shown to preserve every feature of a common observed law, so the displayed range is not claimed as a sharp identified set either.

Negative controls can provide additional falsification information only under their own exclusion and confounding structure. A variable called a negative control by an assistant or a user does not gain those properties by declaration. The present commercial comparison has not completed a product-specific census of negative-control or sensitivity procedures. No failure to locate an option in these selected manuals establishes that a vendor or an extension lacks it. Completing that survey requires matched originals and explicit parameter mappings, not broader interpretation of an existing bootstrap confidence interval.

\subsection{Commercial alternatives and a proposed integration contract}
\label{exec20:med-integration}
Stata supplies a documented parametric causal-mediation workflow, including potential-outcome means, interaction, controlled effects, multiple output scales and several variance options. Mplus supplies a different modeling environment with latent-variable and conventional path capabilities alongside a counterfactual branch. These are substantive strengths against a proposal whose current executable evidence in this comparison consists of exact authored references. There is no measured accuracy, speed or integration advantage over either vendor. A future adapter should preserve their documented scientific conventions before comparing matched targets.

Three further candidate families need matched technical originals: SAS/STAT CAUSALMED, IBM/AMOS and PROCESS. The available evidence does not support a complete edition-specific comparison for these families. The PROCESS author website establishes an independent extension for SPSS, SAS and R, and explains that its documentation is version sensitive. This extension must not be counted as a native IBM or SAS capability without verifying the concrete host, version, terms and implementation. The author's book page is a discovery source; the book and linked scholarly papers are not represented as inspected originals here. Missing technical evidence leaves a comparison unresolved and does not establish product absence.~\citep{exec20_med_process_lead}

The proposed library request should identify treatment and mediator intervention values, the permitted information clock, graph assumptions, natural/controlled/interventional/path-specific target, effect scale, interaction structure and target population. It should also identify resampling units, covariance options, sensitivity parameters and source versions. Outputs should preserve the four nested means when available, compatible decompositions, denominator/sign diagnostics, uncertainty type and unsupported status. A client must not receive a sensitivity range under a confidence-interval field or a latent-response coefficient product under a repayment-risk field.

CLI and MCP transport should preserve that same typed object, including failure and refusal. An agent can compare requested semantics with a vendor command and expose mismatches. It cannot supply missing ignorability, historical entitlement or calibration through a fluent explanation. The first public prototype experiments should check the exact interaction means, the two indistinguishable observed worlds, covariate averaging, conflicting effect signs and sensitivity arithmetic. Subsequent licensed numerical comparisons should measure fitting, conversion, review time and failed requests under frozen versions and generated mechanisms. Bank data and measured commercial value become relevant after those public experiments establish why a prototype deserves a costly institutional trial.
